% Part: first-order-logic
% Chapter: axiomatic-deduction
% Section: rules-and-proofs

\documentclass[../../../include/open-logic-section]{subfiles}

\begin{document}

\iftag{FOL}
      {\olfileid[zh]{fol}{axd}{rul}}
      {\olfileid[zh]{pl}{axd}{rul}}

\olsection{规则与\usetoken{P}{derivation}}

\begin{explain}
  公理化!!{derivation}或许是逻辑学中最简单的!!{derivation}系统。!!^a{derivation}就是一个!!{formula}序列。要算作!!a{derivation}，序列中的每个!!{formula}要么是某条公理的一个特例，要么由推理规则从序列中它之前的一个或多个!!{formula}得出。!!^a{derivation}!!{derive}它的最后一个!!{formula}。
\end{explain}

\begin{defn}[!!^{derivability}]
  若 $\Gamma$ 是 $\Lang L$ 的一个!!{formula}集合，则从~$\Gamma$ 出发的一个
  \emph{!!{derivation}}就是!!{formula}的有限序列 $!A_1$，
  \dots,~$!A_n$，其中对每个 $i \le n$，下列条件之一成立：
\begin{enumerate}
\item $!A_i \in \Gamma$；或
\item $!A_i$ 是一条公理；或
\item $!A_i$ 由推理规则从某个 $!A_j$（和 $!A_k$）推出，其中 $j < i$（以及
  $k < i$）。
\end{enumerate}
\end{defn}

  怎样的!!{derivation}才算正确，取决于我们允许哪些推理规则（当然还有我们视哪些为公理）。而推理规则是一种 if-then 陈述，它告诉我们：在特定条件下，!!a{derivation}中的一步~$A_i$ 是一个正确的推理步骤。

\begin{defn}[推理规则]
  一条\emph{推理规则}给出什么算作从~$\Gamma$ 出发的!!a{derivation}中一个正确推理步骤的充分条件。
\end{defn}

  例如，由于任何满足 $!A \in \Gamma$ 的单元素序列 $!A$ 平凡地就算作!!a{derivation}，下面可能就是一条非常简单的推理规则：
\begin{quote}
  若 $!A \in \Gamma$，则 $!A$ 在从~$\Gamma$ 出发的任何!!{derivation}中总是一个正确的推理步骤。
\end{quote}
  类似地，若 $!A$ 是某条公理，则单独的 $!A$ 就是!!a{derivation}，所以这也是一条推理规则：
\begin{quote}
  若 $!A$ 是一条公理，则 $!A$ 是一个正确的推理步骤。
\end{quote}
  若推理规则诉诸于出现在所考虑步骤之前的!!{formula}，情况就更有意思了。下面的规则叫做\emph{分离规则：}
\begin{quote}
  若 $!B \lif !A$ 与 $!B$ 出现在!!{derivation}更靠前的位置，
  则~$!A$ 是一个正确的推理步骤。
\end{quote}
  如果这是唯一的推理规则，那么我们上面关于!!{derivation}的定义就相当于：$!A_1$, \dots,~$!A_n$ 是!!a{derivation}，当且仅当对每个 $i \le n$，下列条件之一成立：
\begin{enumerate}
\item $!A_i \in \Gamma$；或
\item $!A_i$ 是一条公理；或
\item 对某个 $j < i$，$!A_j$ 是 $!B \lif !A_i$，且对某个 $k < i$，
  $!A_k$ 是~$!B$。
\end{enumerate}
  最后一条说的是：$!A_i$ 由分离规则从~$!A_j$（$!B \lif !A_i$）和 $!A_k$
  （$!B$）推出。若我们能从 $1$ 走到~$n$，且每走一步都找到!!a{formula}~$!A_i$，它要么在~$\Gamma$ 中、是一条公理，要么有一条推理规则告诉我们它是一个正确的推理步骤，那么整个序列就算作一个正确的!!{derivation}。

\begin{defn}[!!^{derivability}]
  若存在从 $\Gamma$ 出发、以 $!A$ 结尾的!!a{derivation}，则!!{formula}~$!A$ 称为从 $\Gamma$ 出发\emph{!!{derivable}}，记作
  $\Gamma \Proves !A$。
\end{defn}

\begin{defn}[定理]
  若存在从空集出发的!!a{derivation}，则!!^a{formula}~$!A$ 是一条\emph{定理}。若 $!A$ 是定理，则记作 $\Proves !A$；若不是则记作 $\Proves/ !A$。
\end{defn}


\end{document}
